% !TeX program = lualatex
% Compile twice: lualatex 177.tex
% All references and prose are embedded; no BibTeX database or external inputs.
\documentclass[11pt,a4paper]{article}
\usepackage[margin=24mm,headheight=14pt]{geometry}
\usepackage{fontspec}
\setmainfont{Latin Modern Roman}
\setsansfont{Latin Modern Sans}
\setmonofont{Latin Modern Mono}
\usepackage{amsmath,amssymb,mathtools}
\usepackage{unicode-math}
\setmathfont{Latin Modern Math}
\usepackage{microtype}
\usepackage{enumitem,xurl,needspace}
\usepackage[dvipsnames]{xcolor}
\usepackage{hyperref}
\hypersetup{unicode=true,colorlinks=true,linkcolor=MidnightBlue,urlcolor=MidnightBlue,
 pdftitle={500 Mathematical Problem Statements},pdfauthor={Source-annotated compilation},bookmarksnumbered=true}
\usepackage{bookmark}
\usepackage{tocloft}
\setlength{\cftsecnumwidth}{3em}
\cftsetpnumwidth{2.5em}
\cftsetrmarg{4em}
\usepackage{titlesec}
\titleformat{\section}{\Large\bfseries\raggedright}{\thesection}{0.7em}{}
\usepackage{fancyhdr}
\pagestyle{fancy}\fancyhf{}
\fancyhead[L]{\small\sffamily Mathematical problem statements}
\fancyhead[R]{\small Source edition 22}
\fancyfoot[C]{\thepage}
\setlength{\parindent}{0pt}
\setlength{\parskip}{5pt plus 1pt}
\setlength{\emergencystretch}{3em}
\setcounter{tocdepth}{1}
\setcounter{secnumdepth}{1}
\setlist[itemize]{leftmargin=3em,itemsep=5pt,topsep=4pt,parsep=0pt}
\allowdisplaybreaks
\newcommand{\N}{\mathbb{N}}
\newcommand{\Z}{\mathbb{Z}}
\newcommand{\Q}{\mathbb{Q}}
\newcommand{\R}{\mathbb{R}}
\newcommand{\C}{\mathbb{C}}
\newcommand{\F}{\mathbb{F}}
\newcommand{\A}{\mathbb{A}}
\newcommand{\PP}{\mathbb P}
\newcommand{\E}{\mathbb{E}}
\newcommand{\eps}{\varepsilon}
\newcommand{\dd}{\,\mathrm d}
\newcommand{\sref}[2]{\hyperlink{source:#1:#2}{[S#2]}}
\newcommand{\eref}[2]{\hyperlink{extra:#1:#2}{[E#2]}}
% Turn the next switch off for a shorter reading copy. The quotations preserve
% every exactTarget field, including qualifications omitted from a short summary.
\newif\ifshowcatalogue
\showcataloguetrue
\newcommand{\cataloguescope}[1]{\ifshowcatalogue
 \paragraph{Exact scope in the supplied catalogue.}
 \begin{quote}\small #1\end{quote}\fi}

% Standalone notebook wrapper; original problem section retained verbatim.
\numberwithin{equation}{section}
\hypersetup{pdftitle={Research attempt -- problem 177},
 pdfauthor={Research notebook; source compilation retained}}
\fancyhead[L]{\small\sffamily Research notebook}
\fancyhead[R]{\small\sffamily Problem 177 | 27 September 2026}
\begin{document}
\setcounter{section}{176}
% ================================================================
% problemId: problem.soliton-resolution-for-the-focusing-energy-critical-wave-equation-without-radial-symmetry
% source releaseRank: 177; source releaseStatus: open_with_solved_subcases
% exactTarget SHA-256: b5c46b052f10a0f386bbef61a890d38c025c70bfdbc23e2acbba24abee9f37d3
\clearpage
\section{\texorpdfstring{Soliton Resolution for the Focusing Energy-Critical Wave Equation without Radial Symmetry}{Soliton Resolution for the Focusing Energy-Critical Wave Equation without Radial Symmetry}}\label{177}
\subsection{Definitions and mathematical statement}
For $d\in\{3,4,5\}$ consider $u_{tt}-\Delta u=|u|^{4/(d-2)}u$ with
$\sup_{t\ge0}\|(u(t),u_t(t))\|_{\dot H^1\times L^2}<\infty$.
A stationary profile $Q\ne0$ solves $-\Delta Q=|Q|^{4/(d-2)}Q$ with finite energy; a Lorentz transform has velocity $\ell$ with $|\ell|<1$. The assertion is a decomposition in the energy norm, for \emph{all} $t\to\infty$, into a free wave, finitely many rescaled and translated Lorentz profiles, and an error tending to zero. The scales and centers satisfy
\[
 \lambda_j/t\to0,\quad c_j/t\to\ell_j,\quad
 \frac{\lambda_j}{\lambda_k}+\frac{\lambda_k}{\lambda_j}
 +\frac{|c_j-c_k|}{\lambda_j}\to\infty\ (j\ne k).
\]
The vector profiles and scaling of their velocity components are exactly those of the cited Case II. Convergence only along selected times is not the target; neither radial symmetry nor ground-state-only profiles are assumed.
\par\smallskip\textit{Formulation and concept references:} \sref{177}{1}.
\cataloguescope{For d in \{3,4,5\}, consider the real focusing wave equation u\_tt-Delta u=|u|\textasciicircum{}(4/(d-2))u on R\textasciicircum{}d, with energy data in dot-H1 x L2 and the usual strong energy solution on [0,infinity), uniformly bounded in that norm. Assertion: every such solution has, for all t->infinity, a decomposition of (u,u\_t) into a finite-energy free wave, finitely many scaled and translated Lorentz transforms of fixed nonzero finite-energy stationary solutions, and a remainder tending to zero in dot-H1 x L2. More precisely, use the global CaseII decomposition(1.10) of DJKM1601.01871v3 with t replacing t\_n: J>=0 is finite, each velocity ell\_j has |ell\_j|<1, lambda\_j(t)>0 with lambda\_j(t)/t->0, c\_j(t)/t->ell\_j, and for j!=k the sum lambda\_j/lambda\_k+lambda\_k/lambda\_j+|c\_j-c\_k|/lambda\_j tends to infinity as in(1.9). No radiality, compact support, ground-state-only or distinct-velocity hypothesis is imposed.}
\subsection{Short English statement}
Must every bounded global energy solution eventually separate into radiation and finitely many noninteracting traveling stationary profiles, with an error vanishing at every late time?
% BEGIN RESEARCH ADDITION ATTEMPT-177
\subsection{Research attempt}

\paragraph{Current frontier.}
Duyckaerts--Jia--Kenig--Merle establish the specified decomposition along a
sequence, with the global Case II vector profiles in (1.10)
\sref{177}{1}. This source has been inspected with its scale, center,
velocity, and orthogonality conventions intact. A 2026 result of
Collot--Duyckaerts--Kenig--Merle treats compactly supported data under
additional assumptions that the sequential profiles are ground-state
traveling waves with distinct velocities. It obtains a dichotomy and rules
out the collision alternative when there are at most two solitons
\eref{177}{1}. Those restrictions are specifically absent here. The
attempt targets the passage from selected times to every late time.

\paragraph{Attack route.}
A channel-of-energy route must rule out excursions between concentration
profiles. A modulation route needs rigidity and control near arbitrary
stationary states, including unstable directions and possible interactions.
A third route combines a sequential decomposition with a quantitative
no-return or finite-excursion principle. We formulate the latter precisely,
prove a sufficient analytic criterion, and test whether ordinary continuity
and an apparently finite dissipation integral would suffice.

\paragraph{Attempt: a defect that includes all asymptotic conditions.}
Fix the free radiation $u^L$, the finite list of nonzero stationary profiles,
and the velocities $\ell_j$ supplied along one sequence by Case II.
For a profile write, with the source's vector convention,
\[
 \vec Q_{j,\lambda,c}(x)=
 \left(\lambda^{1-d/2}Q_{\ell_j}((x-c)/\lambda,0),
       \lambda^{-d/2}\partial_tQ_{\ell_j}((x-c)/\lambda,0)\right).
\]
For $t\ge1$ and parameters $\lambda_j>0$, $c_j\in\R^d$, define
\begin{align*}
 \mathfrak d(t;\lambda,c)
 &=\left\|\vec u(t)-\vec u^L(t)-\sum_{j=1}^J
                              \vec Q_{j,\lambda_j,c_j}\right\|_{\dot H^1\times L^2}\\
 &\quad+\sum_j\left(\frac{\lambda_j}{t}
                         +\left|\frac{c_j}{t}-\ell_j\right|\right)\\
 &\quad+\sum_{j<k}
 \left(\frac{\lambda_j}{\lambda_k}+\frac{\lambda_k}{\lambda_j}
                   +\frac{|c_j-c_k|}{\lambda_j}\right)^{-1},
 \qquad D(t)=\min\{1,\inf_{\lambda,c}\mathfrak d(t;\lambda,c)\}.
\end{align*}
For $J=0$, all sums and the parameter infimum disappear. The sequential
result gives $D(t_n)\to0$ for a suitable sequence $t_n\to\infty$.
For this fixed list of profiles, the asserted all-time decomposition would
follow from $D(t)\to0$. Indeed choose parameters within $1/t$ of the
infimum; every nonnegative term then tends to zero. No minimizer or compact
parameter space is assumed. Conversely, an all-time decomposition using
this list forces $D(t)\to0$.

The distance is measurable: parameter space is separable and the profile
maps are continuous in energy norm, so the infimum can be taken over a
countable dense set. This construction includes the actual scale and
center conditions; an energy-norm distance alone could overlook them.
It is a sufficient approach with the list fixed by one sequence, not a
proof that arbitrary sequential lists must already coincide.

\paragraph{Attempt: a finite-excursion lemma.}
Suppose that $s\mapsto D(e^s)$ is uniformly continuous on $[0,\infty)$
and
\begin{equation}\label{eq177-integral}
                         \int_1^\infty D(t)^2\,\frac{dt}{t}<\infty.
\end{equation}
Then $D(t)\to0$.

\textit{Proof.} Put $d(s)=D(e^s)$. If $d$ does not tend to zero, there
are $\eps>0$ and $s_n\to\infty$ with $d(s_n)\ge\eps$.
Uniform continuity provides a fixed $\delta>0$ on which
$d(s)\ge\eps/2$ whenever $|s-s_n|\le\delta$. Pass to a subsequence with
disjoint intervals. Their contributions to $\int_0^\infty d(s)^2ds$
are each at least $\delta\eps^2/2$, contradicting
\eqref{eq177-integral}. \hfill$\square$

More generally, it suffices that each excursion $D(t)\ge\eps$ forces
$D\ge\eps/2$ on an interval of uniformly positive logarithmic length.
This weaker persistence property could be the actual PDE input. The lemma
then gives the same conclusion without assuming uniform continuity at all
values of $D$. It is crucial that the length be measured in the measure
appearing in \eqref{eq177-integral}.

\paragraph{Attempt: an explicit false upgrade.}
Ordinary uniform continuity in $t$ cannot replace logarithmic uniform
continuity. Let $\psi\in C_c^\infty((-1/4,1/4))$ satisfy
$0\le\psi\le1$ and $\psi(0)=1$, and put
\[
                    D_0(t)=\sum_{n\ge2}\psi(t-2^n).
\]
The supports are disjoint and locally finite. Thus $D_0$ is bounded and
uniformly Lipschitz in ordinary time. It vanishes between pulses, so it has
a sequence tending to zero, and
\[
 \int_1^\infty D_0(t)^2\frac{dt}{t}
       \le C_\psi\sum_{n\ge2}2^{-n}<\infty.
\]
Nevertheless $D_0(2^n)=1$ for every $n$. On logarithmic time the pulses
have widths comparable to $2^{-n}$, so the required uniform continuity
fails. This is a rigorous counterexample to the proposed analytic upgrade,
not a constructed wave solution.

For the PDE, local stability is naturally tied to a concentration scale
$\lambda(t)$, rather than to $t$. Where $\lambda(t)/t\to0$, an interval
of physical length $O(\lambda(t))$ has logarithmic length
$O(\lambda(t)/t)$, which tends to zero. Thus even a uniform local theorem
at the natural soliton scale need not supply the persistence required by
\eqref{eq177-integral}. Changing to an intrinsic clock such as
$\int dt/\lambda(t)$ requires a corresponding finite-dissipation estimate;
one cannot change the clock while retaining the same integral without proof.

\paragraph{Attempt: a no-return alternative.}
A second sufficient input, closer to modulation arguments, would be:
for every $\eps>0$ there are $\delta>0$ and $T<\infty$ such that
\begin{equation}\label{eq177-noreturn}
 s\ge T,\ D(s)<\delta\quad\Longrightarrow\quad
                        \sup_{t\ge s}D(t)<\eps.
\end{equation}
The known sequence provides an $s=t_n$ satisfying the antecedent, so this
would prove $D(t)\to0$ immediately. It is an additional dynamical theorem,
not a compactness consequence. Conserved energy by itself gives no positive
quantity equal to $D^2$ whose total integral is finite, and unstable
stationary states make a forward-invariant small tube especially delicate.

The 2026 compact-support result \eref{177}{1} illustrates the relevance of
collisions and no-return analysis, but it cannot be invoked for arbitrary
sign-changing stationary states, equal velocities, or an unrestricted
number of profiles. None of those cases is removed from the definition of
$D$ above. No general proof of \eqref{eq177-integral}, its persistence
hypothesis, or \eqref{eq177-noreturn} has been obtained here.

\paragraph{Stress test.}
The argument preserves the velocity-component scaling $\lambda^{-d/2}$,
not the scaling of the position component. Orthogonality is encoded by the
reciprocal of the exact divergent expression; the inverse tends to zero
precisely when that expression diverges. Parameters need not attain the
infimum and need not be continuous for the target as printed.

Neither a subsequence of small defects nor conservation of the energy norm
for a linear model proves the nonlinear all-time assertion. The pulse
example shows a failure even under an extra finite weighted integral.
The analytic lemma is valid, but the hypotheses needed to apply it have
not been inferred from the equation. No radiality or ground-state-only
reduction is being made.

\paragraph{Outcome.}
\textbf{Outcome: Conditional reduction.}

The attempt constructs an exact asymptotic defect and proves two sufficient
upgrade criteria: finite weighted excursions with logarithmic persistence,
or a quantitative no-return property. It supplies a counterexample to a
weaker continuity-based upgrade. The unresolved work is PDE control of
excursions for the full nonradial profile class. These deductions organize
a possible proof route; they are not a new soliton-resolution theorem or
a claim of priority for the underlying analytic lemma.
% END RESEARCH ADDITION ATTEMPT-177
\subsection{Sources}
\begin{itemize}\raggedright
\item[\textnormal{[S1]}]\hypertarget{source:177:1}{} Soliton resolution along a sequence of times for the focusing energy critical wave equation,1601.01871v3,7June2017.
\url{https://arxiv.org/abs/1601.01871}.
{\small\textit{Catalogue formulation source.}}
% BEGIN RESEARCH ADDITION SOURCES-177
\item[\textnormal{[E1]}]\hypertarget{extra:177:1}{} Charles Collot,
Thomas Duyckaerts, Carlos Kenig, and Frank Merle,
\emph{Non-radial soliton resolution for the energy critical wave equation
for initial data with compact support}, Discrete and Continuous Dynamical
Systems 56 (2026), 501--556; published online 8 June 2026,
abstract and stated ground-state/distinct-velocity dichotomy.
\url{https://www.aimsciences.org/article/doi/10.3934/dcds.2026115}.
{\small\textit{Consulted 27 September 2026. Compact support, the profile
restrictions, and the at-most-two condition for excluding collisions are
not dropped.}}
% END RESEARCH ADDITION SOURCES-177
\end{itemize}

\end{document}
